\chapter{The Importance of SCEPTICISM}\label{ch:language}
\markboth{The Importance of Scepticism}{The Importance of Scepticism}

\chapquote{``I like to regard myself as someone who's capable of critical thought, that is to say, who can evaluate claims.''}{Bill Nye}

On 30~January 2008 the History Channel put out a programme that, entirely by accident, taught this
subject more about method than the previous fifty years of literature had managed. It was
called \emph{UFO~Hunters}. It ran for three seasons and thirty-nine episodes, it was narrated by
James Lurie, it carried the tagline ``Hoax or History?'', and it was cancelled in the spring of
2009 with four episodes still unaired. By the usual
standards of the genre it should have been forgettable. What made it matter was the casting.

Instead of assembling a panel of enthusiasts who agreed with one another --- the format that has
since become the industry default --- the producers built the team out of three men who could not
possibly agree. \hi{One believer. One sceptic. One who had not made up his mind.} The believer was
William~J.\ Birnes, the show's co-creator and the publisher of \emph{UFO~Magazine}, a man who
arrived at every location already convinced that something was there. The sceptic was
Ted~Acworth, a mechanical engineer trained at MIT, who arrived convinced that something
conventional was there and brought instruments to prove it. Between them stood Pat~Uskert, a
researcher and diver who had come to the subject through his own unresolved sighting and had spent
the years since being unable to settle the question either way. Fans called them the Dynamic
Trio. What they actually were is the cheapest working model of scientific method ever put on
basic cable.

Understand what that structure does. Birnes cannot quietly drop the mundane explanation, because
Acworth is standing next to him with a spectrometer. Acworth cannot wave a case away as
``probably a balloon'', because Birnes will ask him which balloon, launched by whom, from where.
And neither of them can win by force of personality, because the man in the middle has no side to
defend and no reason to be impressed. Uskert is the reader. He is the only one of the three whose
mind is available to be changed, which makes him the only one whose conclusion carries any
information at all. Every fallacy dissected in this chapter is a fallacy that one of those three
men was structurally prevented from committing --- not because he was virtuous, but because
somebody standing beside him had an incentive to catch him at it. That is the entire mechanism of
science, reduced to three folding chairs and a Winnebago.

The pedigree is worth noting because it explains why the show could do this and its successors
cannot. It grew out of a 2006 History special, \emph{Deep Sea UFOs}, in which Birnes and Uskert
were sent to look at underwater sighting claims; it was the first cable reality series shot
entirely in full high definition and native widescreen, executive-produced by Jon~Alon Walz. The
field team changed over its run --- the biologist Jeff~Tomlinson worked the first season,
Acworth left after the second, and Kevin~Cook took the technical seat for the third --- which
turns out to be the most instructive thing about it. \hi{When the sceptic's chair emptied, the
programme's evidential standard fell with it.} You can watch the quality of the reasoning
degrade in real time across thirty-nine episodes, and that degradation is a controlled experiment
nobody intended to run. Birnes has claimed the series was killed for going too far, pointing to
the Dulce Base episode; the mundane explanation is that the network cut it on 20~May 2009 and
burned off the remaining episodes on 29~October. I have no way to settle which is true, and the
fact that I cannot is itself the first lesson of this chapter.

What History did with the format afterwards is the other half of the story. The network kept the
subject and threw away the argument. \emph{Ancient Aliens} premiered the following spring with no
sceptic in the cast, no instruments, and no possibility of a negative result, and it has now run
longer than \emph{UFO~Hunters} lasted by a factor of five. Two decades of History Channel
programming grew out of that fork in the road, and the branch that survived was the one with the
dissenter removed. That is worth sitting with, because it tells you exactly what the market pays
for and exactly what it does not.

So I am going to use the three of them as instruments all the way through this chapter. Acworth
will hold the burden of proof and the demand for falsifiability. Birnes will hold confirmation
bias, because he had it magnificently and sincerely, which is the only interesting kind. Uskert
will hold the difference between doubt that investigates and doubt that merely refuses. Keep all
three in the room while you read. If you can only hear one voice in your head when you assess a
claim, it is almost certainly your own.

% ---------------------------------------------------------------
\section{The Fallacy of Language}\label{sec:fallacylanguage}

Start with the question the trio could never ask each other, because the words would not carry it.
Put Birnes, Acworth and Uskert in a room, ask them ``do you believe in UFOs?'', and you will get
three answers that agree completely on the facts and appear to be a violent disagreement. All
three men had seen unidentified things, or had read the files of people who had; all three agreed
the reports existed and that some of them were unresolved. What they disagreed about was what the
question was secretly asking. Birnes heard ``is something extraordinary going on?'' and said yes.
Acworth heard ``are these alien spacecraft?'' and said no. Uskert heard the question as actually
phrased and could not answer it. That is not three positions. That is one set of facts run
through three different translations, and the translation happened before anybody spoke.

For all the amazing, wonderful things that human language has made possible.
Its ability to convey meaning and emotion. Its capacity to transmit culture and
history through writing and literature. Human language is fundamentally flawed.
Flawed in its ability/hindrance to inject implicit meaning into otherwise
explicit questions. For example, in a 1995 survey of unknown sample size
conducted by Scripps Howard News Service and Ohio State University;\footnote{See
the References section, entry~1.} the populous surveyed was presented with the
following question:

\pullquote{``Do you believe in UFOs?''\\ \normalfont Yes \quad No}

The results suggested an even split of 50\% with 34\% of Seniors (Age 65+),
and 45\% of Church Attendees answering in the affirmative. Such a question when
posted in the following format creates many problems. Most often the phrase
``Do you believe in UFOs?'' is interpreted to read ``Do you believe in UFOs AND
that UFOs are alien spacecraft that have come to planet earth from a distant
star system in effort to take over the world?'' In which case the majority would
answer HELL NO! Such a fallacy of language is a common problem when conducting
public surveys on any number of topics. Much in the way that questions of ghosts
or Bigfoot will automatically conjure explanations for such phenomena as the
restless undead or hairy woodland beasts who make their living by ripping their
prey in half limb by limb. A better written, and more informative question for
example should have read:

\pullquote{``Do you believe UFOs are ground in reality?''\\ \normalfont Yes \quad No}

In which case anyone proclaiming that UFOs are nothing more than trickery of the
mind or pranks conducted by irresponsible hoaxers would still be obliged to
answer ``Yes'' to the question posed, because both explanations are ground in
one's rational understanding of reality. It's only when the idea that UFOs are
spaceships with little green men from some distant star system do people tend to
lean in the side of non-belief. When in fact explanations for the existence of
UFOs are the least interesting things that can be learned from such questions.
It doesn't matter. It does not matter whether UFOs are result of government
funding top-secret black budget projects, or aliens from a distant galaxy, time
travelers from the year 2939, or simply hallucinations of the mind. In each case
regardless of whether the explanation is the one true and genuine answer for the
existence of UFO's, all will open new areas of study and investigation. If it's the government and men in black, then ok what
sort of social or political implications do such an apophony pose? If UFOs are
REALLY time travellers from the future, then what can be learned about the physics
of today that will lead us into the time travel of tomorrow? If UFOs are nothing
more than flawed human perception, what does that say about ourselves and who we
are as a species? These are the types of questions that the field of Ufology is
truly interested in.

Another problem with the question posed from the survey is the constraint of
binary answers. Forcing a person to answer in the affirmative 100\% Yes or
100\% No without allowing any grey area in between. These grey areas that are so
important for such provocative questions. Like asking for instance: Are You
Religious? Are You Liberal? or Are You Gay? Quick! Answer now yes or no. It's
complicated. And complicated questions require careful consideration, thought,
and further reasoning. Difficult questions which allow only binary answers are
fundamentally flawed; in the way human language is fundamentally flawed and
should not be taken as anything more than mere polls for the public enjoyment.
Any quoted ``scientific'' poll pertaining to society's belief in the paranormal
should contain prior planning that would consider such limitations as the fallacy
of language to induce otherwise unstated meaning. And to consider other aspects
like demographics that could potentially play with the results and induce a bias
into said conducted survey.

I hesitate in making assertive assumptions or to claim that any of the surveys
found in the sourced reference material are in some way flawed. The referenced
material which was published in January 2000 by The Mutual UFO Network examines
the previous 50+ years of publicly conducted polls and surveys on the topic of
UFOs. Polls which were conducted both by universities and by private corporations
such as journals and magazines. I am hesitant to dismiss any or all these surveys
as quote ``scientific'' without doing further research into each one. I am unable
and without the resources required to personally conduct any survey of my own.
Unable to test whether the wording of questions or the number of available options
for response in fact play significant role in the information obtained. But I
strongly encourage the consideration of such arguments I have made when overseeing
the formation of any publicly conducted survey, whether done for scientific
research or mere fun and games.

The fallacy of language plays a huge role in all of science and how we go about
inferring results and interpreting data from studies, and should always be taken
into careful consideration.

Watch what \emph{UFO~Hunters} did about this, because it did something right almost by
instinct. The team did not ask witnesses whether they had seen a UFO. They asked what direction
the thing came from, how long it was in view, what it passed in front of, what the sound was like,
whether anything on the property stopped working. Acworth then went and measured whatever could be
measured. None of those questions contain a hypothesis, which is precisely why they are worth
asking, and it is why the interview protocol in Chapter~\ref{ch:ufocsi} looks the way it does.
\hi{A question with a theory hidden inside it will return that theory to you, dressed as a
finding.}

% ---------------------------------------------------------------
\section{The Burden of Proof}\label{sec:burdenproof}

Ted Acworth is the reason this section is short and hard. He had, in the phrase the show never
used, a ``testing habit'': presented with a claim, his first move was to ask what physical measurement
would distinguish it from the boring alternative, and his second was to go and take that
measurement. He did not argue people out of their sightings. He built a rig, ran it, and reported
the number. Sometimes the number went against him and he said so on camera, which is a great deal
more than the field's professional debunkers have ever managed.

When it comes to science; theories are the gold standard in evidence proof, and
the extent of human knowledge. Yes, I know what you are thinking, and I completely
agree. They could not have chosen a more misleading term to represent such a thing.
Through the course of everyday conversation, the average individual will use and
interpret ``theory'' to mean something that a person only thinks is true. Or maybe
there is evidence it may be true, but no actual proof that it is. In other words:
mere speculation. In science, this is what we call a hypothesis. Scientific
hypotheses are founded on the notion that the scientist has reason to assume or
believe that certain outcomes will result from the conduction of his experiment.
But until they actually do the experiment, their hypothesis will remain unproven.

In the field of Ufology there are NO theories. None. There are only rumors,
speculation, and any number of explanations for the phenomena. These are the UFO
hypotheses; the ancient alien hypothesis, the abduction hypothesis, and the
government classified top secret technology hypothesis, the reptilian hybrid
hypothesis of David Icke to name a few. \hi{This is
because of one simple fact that the entire foundation of science rests upon:
Science MUST be falsifiable.} Any experiment in which the scientist proposes one
outcome will happen, must allow for the chance and possibility (even if very
small) that the entire opposite will happen. This is the scientific premise
Ufology lacks, and the one thing preventing it from become an officially
recognized field of inquiry. The ability to test and form falsifiable experiments on the hypotheses
proposed as possible explanations for the subject matter. That is why the burden
of proof lies on the Ufologist who studies this subject. Much in the manner that
Hitchens's Razor has inspired many variations on the following quote:
``Extraordinary claims require extraordinary evidence.''

\begin{funfact}
A philosophical razor is not in any way sharp or pointy but is simply an idea or
rule of thumb to abide by when drawing conclusions using a skeptical mind.
Occam's Razor for instance is likely the best-known example of such. Occam's
razor states that in the case of 2 possible solutions or answers to a question;
chances are the simpler of the two is more likely to be right. Occam's razor
however is often used as a \term{strawman} argument to simply debunk UFO's by
asserting that conventional wisdom can explain it away. Swamp gas, weather
balloons, the planet Venus, hallucinations, mass hallucinations, hysteria,
hoaxes, airplanes, ball lightning, frisbees and flares you name it; it's all
rubbish!
\end{funfact}

Take the case of the hypothetical husband for instance whose fingerprints are
found on the weapon used to kill his wife. According to Occam's Razor, that
man is guilty of murder! How else could one explain the existence of his
fingerprints on the weapon OTHER than the simplest answer: he killed his wife.
Preposterous you say! That man loved his wife. You suggest he was framed by the
true murderer, for what other explanation could there be for such an egregious
act from a loving husband? The argument you now present in defense of the
husband's innocence entirely contradicts the notion first proposed under Occam's
razor as you've introduced new variables into the equation. Yet how easy it would
be to now assert that this is the NEW simplest answer. Regardless of what any
investigation discovers about the wife's murder, you will ALWAYS be able to
assert that the new explanation is now the truth, not because it is simplistic,
but because it is true. The investigation into the murder may in fact discover a
plot that spans 10 years and includes 3 assassins, 2 lovers, 12 kitchen knives,
and a rowboat! Just because the answer is complex does not mean it is false.

However, that does not allow for legitimate argument into the completely opposite
and equally absurd explanation for UFO sightings. Which states that Zoltar from
the planet Koolon has come 14 million miles to Earth to enslave the human race but
has chosen to do so by giving birth to a race of reptilian shapeshifters who have
now infiltrated the highest ranks of the US government and plot our very demise
through the use of unaccountable tax returns, burning more fossil fuels,
and \dots\ pee-pee tapes. No. If we are to find reasonable legitimacy for the
further scientific investigation of UFOs then there must be some common ground
that lies both within the realm of possibility and within the validity of such
explanations not to assert that UFO eye witnesses are crazy drunks.

\hi{It should be understood that the default position of any claim is to reject it
until sufficient proof has been presented for acceptance.} This is known as
the \term{null hypothesis} --- or, in less formal language, the default sceptical position.
This position has no bearing on whether the claim is of positive or negative
declaration. If the claim is that UFOs exist, then the default position is to
reject this belief. This does not mean however you assert that UFOs do NOT exist.
Conversely if the claim is that UFOs are total bogus, then the strategy holds; you
should reject this claim. Similarly, that does not mean you now assert the opposite
is true. \hi{The rejection of a claim is not coequal to the assertion of its
counterpart.}

That is why the burden of proof lies on those who investigate the subject matter.
Because nothing should be assumed or believed, and even more so acted on or
regulated without substantial reasons for doing so. This is the plight of the
Ufologist and we must agree to take it as such. It should also be noted that when
investigation into the subject is conducted it should be done without bias or
motive to either prove or disprove such things. Richard Dawkins has
said---whether he was the first to say or not---``By all means let's be
open-minded, but not so open-minded that our brains drop out.''

% ---------------------------------------------------------------
\section{All the Evidence But No Proof}\label{sec:allevidencebutnoproof}

It should also be noted, that scientific theories; although the gold standard in
evidence: are NOT 100\% bulletproof. That is because any scientific theory that
has gone through thousands of experiments, PhD level research, and any number of
well documented case studies is fundamentally flawed. Much in the same way
language is fundamentally flawed, and for all the same reasons. Scientific
theories are based on human perception and understanding of the world, and human
perception is often misleading. We aren't perfect beings, and neither are our
theories and ideas perfect philosophies.

Even the greatest thoughts from the greatest minds, are merely
good-enough-for-now explanations of the world. For example: Einstein's work is
incomplete. Let me say that a bit louder for the people in the back: Einstein's
work is incomplete! No, I did not say Einstein was wrong. I merely observe that
Einstein was flawed because Einstein was human, and humans are flawed. This is
the very reason there exists 2 distinct and disjoint theories for relativity.
They are: general and special relativity. Both equally as true, and yet only
applicable to the situations each describes best. Depending on what you believe;
Einstein likely went to his grave unable to solve his greatest challenge: The
Theory of Everything. One singular equation that will bind the whole of existence
together. My father for instance does not hold such notions. My father believes
that Einstein fully comprehended within the realm of his consciousness the key to
the universe. But was unable to formulate any equation or combination of language
to describe the things he understood; once again resulting from the imperfections
of these human inventions.

What then can be known if anything at all? That is the question explored by
17th century philosopher Ren\'e Descartes. At a time when so many other
philosophers were making appeals to God, Descartes took an alternative route by
appealing to rationality. He imagined a world fashioned by an evil madman whose
sole purpose for existence was to deceive the inhabitances of his false reality.
This thinking of false existence has been used many times throughout some of the
greatest mind-bending films preying on our fear of the unknown. There is no way
however to either prove or disprove such a theory: it's unfalsifiable.

This thought experiment although raises an even bigger question. If everything we
know is fake from our past, to our loved ones, to our income and fears; then does
any of it even matter? This thought troubled Descartes as he was unable to find a
solution. Eventually he would come to realise that the only thing he can know for
certain in taking doubt of the world around him was that he knew it. In other
words: ``I think therefore I am''---Ren\'e Descartes (1637).\footnote{See the
References section, entry~2.}

This method of thinking would eventually lead to the 21st century version of the
``evil genius''. The brain in a vat connected to a supercomputer. Many great
scientists even think it is possible that we are living in a computer simulation.
Statistically speaking, if there is 1 true genuine universe with intelligent
thinking beings. Then it's possible these beings have reached a stage of
technological development in which they can run simulations of their own reality.
Hundreds maybe even thousands of reality simulations could be running all at once
or in parallel. The chances of any being within each simulation being a resident
of the real genuine universe is simply 1 over the number of total simulations
being run. And so statistically speaking we are far more likely to be a resident
of the false simulated reality then someone who resides in the real world. This is
an idea that will be explored further in Chapter~\ref{ch:dimensions} when I talk
about the science of alternate dimensions.

Again, this idea will raise the question: Does reality even matter if we can't
tell it's fake? The answer to this is surprisingly more simplistic then you may
initially think, and although the meaning of one's own reality is entirely a
subjective predisposition; attempts towards a logically correct answer have led to
new insight. The argument runs like this. If the simulation is perfect, then every test you
could perform on it returns the same answer it would return in a base reality, which means the
distinction has no observable consequences --- and a distinction with no observable consequences is
not a scientific claim, it is a mood. If on the other hand the simulation is imperfect, then the
imperfections are physical, they are measurable, and hunting for them is ordinary science rather
than philosophy. Either the question is untestable and therefore idle, or it is testable and
therefore no longer a question about simulation at all but about the behaviour of matter. Notice
that this is precisely the structure we will be applying to UFO claims for the next four hundred
pages. The gravity you feel, the burn you get from a hot stove, and the grief you feel at a funeral
are identical under either hypothesis. All that matters is
consistency.

I would like to point out that this is also the number one problem that arises
when taking a sceptical approach to thinking. The fact that scepticism is
limitless. If you can doubt your very existence, then most certainly the doubt of
having seen a UFO is far more within the bounds of held belief. So therefore, it
becomes easier to simply assert that the experience was a trick of one's own
senses and move on with your life. It takes real tenacity to stand up against the
in-crowd and proclaim there is more to the situation then meets the eye. And so,
these silent heroes and seekers of the truth should NOT be met with jeering
laughter at the gates of discovery, but instead with praises for the courage to
speak their mind and demand the truth!

What then can be made of the truth in a world distraught by political scandal,
coverup and fake news? All factors are bound tightly together and play significant
roles in the field of science, regardless of the specifics on that field of
inquiry. Whether you choose to believe it or not, the world is round. That is what
scientists call an objective fact. It's true whether you choose to hold that
belief or not. We all know however that there are people on this big blue sphere
who would claim otherwise. The motivation behind every flat earther is unique to
each and should be respected on the grounds of freedom of speech. For it is
freedom of speech that leads to discovery. If citizens were not free to express
new ideas regardless of how detestable they are, then we may all still think the
earth is the centre of the universe.

Consider for instance the case of the hypothetical scientist who honestly thinks
the earth is flat. While there is plenty of proof that it is not, the amount of
given data does not satisfy this individual's threshold of proof. If this
individual holds a higher standard of proof than the rest of the population; there
is no logical reason to be angry at them for holding such beliefs. ASSUMING
however, and I want to make this very, very clear! They are willing to undergo
experiments and test their belief even knowing there is a chance they could be
wrong. That is the divide between the hypothetical scientist and real-world
examples of flat-earthers. The scientist will undergo experimentation, while the
whole of flat earthers resort to the side of cult like assertion of false facts.
They are not willing to experiment, and devise means to test their ideas, they
simply want you to agree with them because.

Now in defence of the hypothetical flat-earther scientist, and anyone else for
that matter holding wildly absurd ideas and a means of testing them, consider the
following:

The scientist believes the earth is flat and has a means of testing it. His
hypothesis is the experiment will result in data to support the flat earth theory.
His experiment has never been done before and is so outside the box it has the
scientific community reeling in anticipation. For if his experiment was inside the
box or somehow conventional, then it would have been done already in which case
this scientist would have the data they need to believe the earth is round.
Currently, they do not.

And so, they undergo said experiment, and after all the data is gathered and the
numbers are crunched they come to the stunning realization they were wrong.
``Ahhha''! The scientist shouts in avid proclamation! ``I was wrong after all''.
``But wait?'' ``What is that??'' The scientist then notices something in the
charts he had not before. Some of data ENTIRELY unrelated to the shape of earth,
and yet it was gathered and clearly visible in the chart. And so, they investigate
the abnormality further and spend the next 10 years of their life founding a never
before seen branch of science based on an abnormality that was only seen as result
of their experiment being so unconventional.

Many of the greatest scientific discoveries of all time resulted as
complete accident, or because the scientists trying to prove one thing discover
something else entirely unrelated. That is the conjecture I put forth. That it is
important for people such as the flat-earther scientist to have a platform on
which they can present new ideas. Because even though they began with a false
belief, they have now found evidence in the contrary THAT EXCEEDS their own
standard of believing. In addition, the rest of the scientific community will
benefit from new insight into totally unexplored territory.

How much evidence though is required to prove things? What does proof even mean,
and how do the beliefs we hold as true impact our behaviour? These are all
interesting questions and dive much deeper into the heart of reality when explored
to their fullest potential. If you have ever had the burden of learning math
proofs in post-secondary, then you will understand exactly what is meant by this.
For everyone else, all that is needed to be known is that math proofs are hard as
(really bad word). For the simple reason that in order to prove anything at all you must first agree
on what you are permitted to assume, and mathematics is the only field honest enough to write those
permissions down at the top of the page. Consider $\pi$. Every schoolchild is told it is the ratio
of a circle's circumference to its diameter, which is true and also not a proof of anything. To show
that this ratio is the same for every circle, that it is irrational, and that it is transcendental
--- meaning it is not the root of any polynomial with integer coefficients --- took mathematicians
until 1882, and Lindemann's proof of the last of those does not involve circles in any recognisable
way. It proceeds instead through the exponential function and a result about algebraic numbers, and
when you finish reading it you have been given no picture in your head whatsoever. You have been
compelled. That is what a proof is: a chain of steps each of which you cannot refuse, arriving
somewhere your intuition never went. The proof for $\pi$ is
entirely logical yet makes no sense.\footnote{See the References section,
entry~3.} Hold on to that discomfort, because it is the standard the word ``proof'' carries, and in
this field the word gets used to mean a blurry photograph and a sincere man.

We must accept this simple fact if we are to move forward in our field of
discovery. The fact that to date there is NO proof for the existence of UFOs or
extraterrestrial beings visiting the planet, but that DOES NOT in the slightest
mean there is not evidence. That is the goal of further chapters within the book.
To explore what evidence has been presented and in an unbiased fashion examine the
possible ways in which we can test such evidence. For then, and only then we may
come to a better understanding of the mysterious world around us.

% ---------------------------------------------------------------
\section{Claims Made from a Slippery Slope}\label{sec:claimsmadeslipperyslope}

The concept of a slippery slope essentially conditions that once endeavour down a
certain path of thought or logical reasoning has begun, it's only a matter of time
before things get carried away. My case in point is the infamous History Channel
TV show: \emph{Ancient Aliens}. Well written and visually stunning; the show
offers little to no actual digestible information. It is done this way on purpose.
From a marketing perspective it does its job, and it does this very well for that
matter, by capturing and captivating audiences with its wildly absurd claims and
ideas. Unfortunately, very little to no actual rigorous scientific experimentation
is done throughout the now 13 seasons totalling over 140 episodes.

Functionally speaking, the show suffers from drastic over simplification, and the
incessant need to blame every unknown mystery of the world on extraterrestrial
intervention. This web of connections the show suggests regarding all the earthly
mysteries is both its greatest success and the worst possible thing for genuine
researchers. It is successful because it offers something for everyone. Maybe you
are into native culture and legends of sky people. They got that. Or perhaps
stories of Atlantis deeply affect your existence. Or maybe you yourself are
subject to modern day UFO sightings. Yes, yes, yes and yes, they got all that
covered plus much more and perfectly wrapped into a nicely edited and dramatically
narrated 45 minutes of verbal diarrhea.

This tactic of filming has led to its popular success and obvious financial gain
evident from the fact the show is now 9 years running strong. It offers little
however for actual discovery aside from the possibility that someone watching from
the fence-line who has never formed an opinion on the subject matter may be
inclined to think ``hey this is interesting, I never even considered that.'' Aside
from that possibility however, not much can be gained by watching the show. Worse
is the likelihood that given the sheer amount of information overload, and the fact
the show regards aliens as the solution to every problem; may turn away some people
who would otherwise be inclined to believe, or at least willing to investigate such
unknowns.

This interconnected web of ``discovery'' only serves to dilute the overall content
and cause sheer rejection in the face of laughter by the masses. There is a reason
``Alien Guy'' became a meme easily found on any search engine and providing answers
to all life's questions. It leads to the assertion that once any serious question of
existence is posed, it's only a matter of time before we start asking ourselves:
``Why is the sky blue?'' ``Why did she dump me?'' ``Why was that calculus test so
damn hard??'' And the answer is always undeniably:

This web of interconnected ideas has other unfortunate downsides. Consider the
possibility that 1, or 2, maybe even 10 things that \emph{Ancient Aliens} presents
as PROOF for alien visitation in the past are ACTUALLY CORRECT. There is no way
the average viewer of the show will be able to discern between these factually
correct findings and the other 100+ things that have been thrown into this
concocted mess of conjecture. And so, the solution is to simply reject the whole of
the story as being nothing more than an entertaining idea, and not to be taken
seriously. Therefore, any potential insight into the world is simply lost as no
investigator, reporter, or investigative reporter will touch certain topics with a
50 ft pole.

Which is why scepticism and scientific investigation are the engines that need to
revolutionize this way of thought. To the show's credit, there have been several
scientific tests conducted on certain episodes, but not nearly enough. One such
test that comes to mind is an artefact that was discovered in an Egyptian tomb.
Archaeologists claimed that this artefact modelled the ancient's representation of
birds. Lead show investigator Giorgio A. Tsoukalos however notes that not a single
bird like feature is displayed on this small metallic trinket, and thus reasons
that this artefact must instead be a model of modern-day aircraft. (Somehow). A
replica of this trinket was produced at a larger scale and tested for flight
characteristics.\footnote{See the References section, entry~4.} To basically no
one's surprise, it was deemed aerodynamic. I'm not saying that the test was not
done scientifically. I merely observe that with modern day understanding of flight,
one can very accurately predict whether the model will confine to the structure of
something that is aerodynamic even before the test begins. While testing this
hypothesis is still critically important, it would be very easy to stage such a
``scientific'' test and embed the results within the show in effort to earn brownie
points from the community of sceptics.

\figwrap[r]{0.40}{The so-called ``Saqqara Bird'' --- a sycamore-wood artefact from a
sixth-century BCE tomb, now in the Egyptian Museum in Cairo. It has a vertical
fin rather than a horizontal tailplane, and no trace of a stabiliser.}{https://commons.wikimedia.org/wiki/File:Photo_2-plane_side_view1.jpg}{fig:saqqarabird}
Equally suspicious is the clear lack of any scientific tests conducted on the show
that prove without a doubt something could have been possible WITHOUT alien
intervention. I don't find the so-called aeroplane model found in the tomb to be
particularly convincing evidence to support the hypothesis, but I must admit the
show tries really hard to persuade otherwise. I think \emph{Ancient Aliens} the show
could otherwise benefit overall if they investigated and experimented with ideas
that were perhaps not so easily predictable. And furthermore, if any tests were
conducted on the show that clearly present flaws to their ideas, these tests should
not be edited out. As even the best ideas when investigated will be found to have
many flaws. I will be making a more in-depth analysis of this show, and rather the
ancient alien hypothesis itself, in Section~\ref{sec:aah} of this book, but for now
I think I've said enough.

As for the methodology of sceptical thinking however, I believe that a balanced
approach is needed. An approach that does not immediately reject certain ideas
because they came from a crazy man or a crazy show. An approach however that does
not automatically assume that solutions to all problems are as easy as a singular
blanket statement. Only through unbiased investigation and unfiltered scientific
study can we reach a potential conclusion of the mysterious. Remember the next
time you encounter an episode or idea from the show \emph{Ancient Aliens}: even
broken clocks are correct twice a day.

% ---------------------------------------------------------------
\section{Correlation vs Causation}\label{sec:correlationvscausation}

In the field of statistics there exists a concept often taught in any or all
introductory classes to the subject. \hi{The basis of this concept states that evidence
for correlation is not the same as evidence for causation.} Take 2 or more things
that can be linearly graphed or somehow strung together in a way that may appear
scientific and unbiased, this however does not PROVE that these things are related.
To illustrate this point, consider the graphs included below.

\Needspace*{4.2in}
\begin{center}
\refstepcounter{figph}\label{fig:corrcausetrio}%
\begin{figpanel}[width=0.96\textwidth]
\centering
{\setlength{\fboxsep}{0pt}\setlength{\fboxrule}{0.5pt}%
\fcolorbox{ufogreen}{white}{\adjustimage{max size={0.465\textwidth}{2.6in}}{fg_corrcausea.jpg}}\hfill
\fcolorbox{ufogreen}{white}{\adjustimage{max size={0.465\textwidth}{2.6in}}{fg_pirates.jpg}}}\par\vspace{3pt}
{\RaggedRight\scriptsize\color{steel}\textbf{\textcolor{ufogreenink}{Figure \thefigph.}}
Two ways of making an unrelated pair of numbers look like a mechanism. Left: two series rising
together, which the chart asserts nothing about and the eye explains anyway. Right: the
number of seafaring pirates plotted against global average temperature --- deliberately absurd, and
a better fit than a good deal of published work. Both charts are reproduced large enough to be read,
which is the whole point of putting them here.\par
\vspace{1pt}{\color{midgrey} Left: Hephestus-1964, CC BY-SA 4.0 ---\ Image sourced from:} \srclink{https://commons.wikimedia.org/wiki/File:Correlation_verses_causation_demonstration_chart.svg}\par
{\color{midgrey} Right: LiamVleck, CC BY-SA 3.0 ---\ Image sourced from:} \srclink{https://commons.wikimedia.org/wiki/File:Pirate_Global_Warming_Graph.gif}\par}
\end{figpanel}
\end{center}
\vspace{6pt}

Look at what those two graphs are actually doing. Each one plots two entirely unrelated
quantities, scales the axes so that the wiggles line up, and lets your visual cortex do the rest.
The second is my favourite, because it is a joke that was designed to be a joke: as the number of
seafaring pirates declined over two centuries, global average temperature rose, and the fit is
frankly better than a good deal of published social science. Nobody thinks pirates cooled the
planet. The reason the graph is funny is that the graph is convincing, and those two facts sitting
side by side should make you permanently suspicious of your own eyes. There are three ordinary
explanations for any correlation before you reach for the interesting one --- coincidence, in which
two trends drift together by luck and the coincidence is guaranteed to occur if you test enough
pairs; reversed direction, in which B causes A rather than A causing B; and the lurking third
variable, in which some unmeasured factor drives both. Industrialisation drives the pirates down and
the temperature up, and it does not need to know that either graph exists. \hi{Correlation tells you
where to look. It never tells you what you have found.} In this field the lurking variable is
usually media attention: sightings spike after a film, a broadcast, or a wave of coverage, so any
graph of reports against anything else is really a graph of how much people were thinking about the
subject that year.

\figwrap[r]{0.42}{The Viking~1 frame of 1976 beside the Mars Global Surveyor image of 2001 --- the same landform, better resolution.}{https://commons.wikimedia.org/wiki/File:Martian_face_viking_cropped.jpg}{fig:facemars}
The Face on Mars which was first photographed in 1976 by the Viking 1 orbiter was
later proven to be false through the advent of technological development that
allowed for improved photographs of the Cydonia region of Mars (see Figure~\ref{fig:facemars}). These improved
photographs clearly showed that the face was indeed nothing more than a small
mountain range when under the right conditions of shadows and lighting will appear
as the face of man carved into the surface of the rock. Anyone who may have
otherwise been on the fence regarding opinion on the subject, and those reporting
that such a thing was real and proof of alien life in our solar system were then
forced into the objective fact that it is not.

Consider however the implications that it was. A detailed face complete with moles,
freckles and even acne scars so perfectly carved into the mountain that it would
have had to been placed there by our space brothers and sisters for us to find.
Would that really be the implication? Or would there still be a conventional way of
explaining the feat? The perfect \emph{correlation} between this hypothetical
carving into the surface of a planet on average some 140 million miles away,
does not lead to the conclusion that the \emph{causation} is aliens. Similar to how
1000 monkeys typing at random on a keyboard for the duration of eternity will
eventually write the complete works of Shakespeare. No logical person would reason
that these monkeys have the literary capacity and creativity of grasping the
profound influence of William Shakespeare. Neither should it be concluded upon
immediate discovery that we are no longer alone in this universe.

However small and seemingly impossible the chance that random natural events can
cause the illusion of perceived intelligence, it should not be completely ruled out
as potential possibility. The possibility that asteroids impacting the surface,
geological activity, or volcanic eruptions could have produced the image of a face
on the surface of Mars is low; astronomically low, yet what are the chances that
life in this one corner of the cosmos would have evolved into the very thing
depicted on this rock, and yet it has.

Yet even the existence of our own intelligence starts to look less remarkable --- less like the
signature of a designer and more like what sufficiently many iterations of a simple rule will
produce given time --- when you consider just what is capable without even a brain. Conway's Game of Life
is a perfect computer science example of what I mean. The rules of the game are
simple. In a grid of arbitrary size there exists several pixels where any pixel that
is on signifies being ``alive'' and all the pixels that are turned off represent
``death''. Rule 1: any live cell with fewer than 2 live neighbors will die in the
next state of the simulation as result of ``loneliness''. Rule 2: any live cell with
greater than 3 neighbors will die from ``overpopulation''. Rule 3: any dead cell
with 3 alive neighbors will be birthed and thus become alive, and finally all other
cells to which none of the mentioned rules apply will retain their current state.

Although the simulation was intended to be nothing more than a fun thought
experiment, with no practical or philosophical intention as admitted by the creator
himself: John Conway.\footnote{See the References section, entry~5.} This simulation
proves to be quite interesting to watch. Over time the development of certain
organisms can be observed. These organisms can be attributed basic characteristics
and behaviours if you will. These behaviours can be used to further classify the
predicted outcomes of the environment, and fans have even gone as far of naming and
categorizing different pixel creatures of the simulation, and their perceived
behaviours. These creatures have no brain they have no desires or purpose for life,
they are nothing more than points of light on an illuminated computer screen. And
yet there is something deeply intrinsic about their movement and characteristics
that causes us to personify them and find relations to the real world for such
beings.

This personification of computer code is what I find especially fascinating since
from a philosophical standpoint there is not much that distinguishes these pixels on
a screen from something like a small bug for instance. We know from the scientists
and people that study these bugs that they do have brains. These brains are very
instinctual and limited only allowing such bugs to do basically just 2 things. Avoid
death or being eaten and find food. From our view of the world we assume that these
bugs are thinking, but what are the chances that the only thing existing in their
tiny bug brain is a few lines of genetic code that causes us to perceive
intelligence in their behaviour? How big a leap would it be to next assume that we
humans are following a larger set of preprogrammed rules that causes the illusion of
behavioural patterns such as speech, music and literature? Well that's essentially
what DNA is. The genetic code for everything that has ever lived on this planet.

I will be exploring in the next chapter in my subsection about AI the idea that
illusions of intelligence can be caused by complete accident, or at the very least
not explicitly outlined by the creator of such AI. All these are perfect examples of
the concept that correlation does not equal causation. A bug can appear to do
intelligent things even if the rules governing its existence have nothing explicitly
to do with its behaviour at all. Or in the case of computer code and AI deep neural
networks, the same principal applies. When it comes to UFOs I will do my best
through the rest of this book to ration out the cases where causation arguments
could lead to false belief on the premises or causes to such incidents.

% ---------------------------------------------------------------
\section{Confirmation Bias}\label{sec:confirmationbias}

Confirmation bias is the tendency to notice, remember, and give weight to the
evidence that agrees with what you already believe, while quietly discarding
everything that doesn't. It is not lying. That is the part people get wrong. The
person doing it is entirely sincere, which is precisely what makes it the most
dangerous failure mode in this entire field.

Here is how it works in practice. A man believes he saw a craft over his house.
He goes looking, and he finds forty-one other reports from the same county that
year. Forty-one! He does not go looking for the two hundred reports from that
county that were resolved as aircraft, satellites, or Venus, because why would he?
He already knows what he saw. The forty-one are evidence. The two hundred are
noise. He has not falsified a thing; he has simply built a case file out of the
half of the data that agreed with him, and the case file looks devastating.

Now flip it, because this is the part the sceptic crowd never wants to hear. The
debunker does exactly the same thing in the opposite direction. He collects the
hoaxes, the drunks, the lens flares, the weather balloons, and he assembles them
into an equally devastating file proving the whole subject is rubbish. He ignores
the radar-visual cases with multiple independent witnesses and instrument
corroboration, because those are inconvenient, and inconvenient data has a way of
becoming invisible. Both men have confirmation bias. Only one of them gets invited
onto television to explain how rational he is.

Bill Birnes was the show's believer, and I want to be fair to him here, because he was also the
reason it existed. He had the energy, the contacts, the willingness to stand in a field at three in
the morning. He was also, reliably, the man who arrived at the location already knowing what the
answer was going to be. When a witness described lights in formation, Birnes reached for craft;
when a document surfaced, he reached for cover-up. That is confirmation bias operating in a
perfectly sincere and hard-working man, and it is why the format put Acworth beside him rather than
letting him narrate alone.

The mechanism has a few reliable symptoms, and you should learn to feel them in
yourself:

\begin{itemize}[leftmargin=1.4em,itemsep=2pt]
  \item \textbf{Asymmetric scrutiny.} You interrogate the evidence against your
  position for hours and accept the evidence for it in seconds.
  \item \textbf{Moving the goalposts.} A test comes back against you, so you
  discover a flaw in the test that you did not mention before you knew the result.
  \item \textbf{Counting hits, forgetting misses.} The psychic got three things
  right. She got fifty things wrong. Nobody wrote those down.
  \item \textbf{Source laundering.} You trust a claim because it appeared in three
  places, without checking whether all three copied the same original.
\end{itemize}

That last one deserves a moment, because it is endemic in Ufology. A single
unsourced sentence written in a paperback in 1978 gets quoted in a magazine in
1985, which gets cited by a website in 2001, which gets referenced in a documentary
in 2014, and by the time it reaches you it has four independent-looking sources and
zero actual evidence behind it. This is how the ``fact'' that the Roswell debris
sprang back to shape, or that Betty Hill's star map matched Zeta Reticuli, hardened
into received wisdom. Trace every claim back to its first appearance. Every single
time. Most of the time the trail ends somewhere embarrassing.

\begin{funfact}
Confirmation bias was demonstrated experimentally by Peter Wason in the 1960s using
a task so simple it is almost insulting. Subjects were shown the number sequence
2--4--6 and asked to work out the rule behind it by proposing further sequences. The
overwhelming majority proposed 8--10--12, then 20--22--24, then 100--102--104,
confirming their guess of ``add two'' over and over. Hardly anyone tried 3--5--9,
which is the only sort of move that could have shown them they were wrong. The
actual rule was simply ``any three increasing numbers''. People do not naturally
test their beliefs. They audition for them.
\end{funfact}

The cure is not to have no beliefs, which is impossible, and anyone claiming to be
free of bias is telling you something useful about their self-awareness rather than
their methodology. The cure is procedural. Decide what would count as disconfirming
evidence \emph{before} you collect any. Write it down. Write down what result would
make you abandon the hypothesis, and then go looking specifically for that. \hi{If you
cannot name anything that would change your mind, you are not investigating, you
are worshipping.}

One more thing, and it applies to me as much as anyone reading this. The most
comfortable position in Ufology is the middle. ``Maybe some of it is real.'' It
feels balanced and reasonable, and it is also perfectly designed to survive any
possible evidence, which should worry you. A position that cannot lose is a position
that cannot learn. I hold mine loosely and I invite you to hold yours the same way.

% ---------------------------------------------------------------
\section{Doubt Driven Scepticism}\label{sec:doubtdrivenskepticism}

There is a variety of scepticism that has nothing to do with inquiry. I call it
doubt driven scepticism, and it is the counterfeit currency of this field. Real
scepticism is a method: you doubt in order to find out. Doubt driven scepticism is a
posture: you doubt in order to be finished. The two look identical from the outside
and they produce opposite results, which is why so many people who consider
themselves scientific are in fact doing nothing of the kind.

The tell is what happens after the doubt is expressed. A genuine sceptic says ``I
doubt that, and here is the test that would settle it.'' The doubt driven sceptic
says ``I doubt that,'' and then goes to lunch. One of these is the opening move of an
investigation. The other is a conclusion wearing an investigation's clothes.

Pat Uskert occupied the seat that this section is really about --- the undecided one. It is the
least glamorous job on the team and the hardest to do honestly, because the undecided position is
the one most easily faked, as I said at the end of Section~\ref{sec:confirmationbias}. What made
Uskert's version of it real was that he moved. He
came back from some cases more convinced and from others visibly less so, and he said which was
which. An undecided investigator who ends every episode exactly as undecided as he began is not
weighing evidence; he is performing neutrality. Uskert kept score, and that is the difference.

Recall from Section~\ref{sec:burdenproof} that the rejection of a claim is not coequal to the assertion
of its counterpart. Doubt driven scepticism is what happens when a person forgets
this. They begin by properly withholding assent from the claim that UFOs are real,
which is correct procedure, and then somewhere along the way they quietly upgrade
that withholding into the positive claim that UFOs are \emph{not} real. That is a
different claim. It has its own burden of proof, and nobody has met it, and almost
nobody notices they have taken on the debt.

Here is the deeper problem, and it is the one I raised in Section~\ref{sec:allevidencebutnoproof}: scepticism has
no natural floor. There is no point at which doubt exhausts itself and says ``well, I
suppose that's settled.'' You can doubt the photograph. Failing that, doubt the
photographer. Failing that, doubt the film stock, the developer, the chain of
custody, the witness's eyesight, the witness's sobriety, the witness's motives, and
failing all of that you can doubt your own existence and go home undefeated.
Unlimited doubt is not rigour. It is a machine for never being wrong, and it costs
nothing to operate.

\begin{funfact}
The original Greek \emph{skeptikos} meant simply ``inquirer'' --- from
\emph{skeptesthai}, to examine or consider. The word did not mean ``one who
disbelieves''. It meant ``one who looks''. Somewhere in the last two thousand years
we managed to invert the meaning of the word entirely, which given the contents of
Section~\ref{sec:fallacylanguage} should surprise nobody at all.
\end{funfact}

So how do you tell the difference in the wild? I use four questions, and they work on
believers and debunkers alike:

\begin{enumerate}[leftmargin=1.6em,itemsep=3pt]
  \item \textbf{Is the doubt specific?} ``That radar return is unreliable'' is
  useful. ``Radar can be unreliable'' is a fortune cookie.
  \item \textbf{Does the doubt cost anything?} A doubt that requires no work and
  risks no reputation is not a scientific contribution.
  \item \textbf{Is it reciprocal?} Apply the same standard to the mundane explanation.
  If ``weather balloon'' is offered with no balloon, no launch record and no recovery,
  that is an unevidenced claim and it does not get a free pass for being boring.
  \item \textbf{What would end the doubt?} If nothing would, the doubt is not about
  the evidence. It is about the doubter.
\end{enumerate}

That third question is where most of the professional debunking of the last seventy
years falls apart. Project Blue Book, which we will get to in
Section~\ref{sec:bluebook}, closed cases as ``probable aircraft'' and ``possible
balloon'' with a frequency that would be laughed out of any other branch of
government science. A probable is not a finding. It is a hypothesis that was never
tested, filed as though it had been, and the sheer administrative confidence of the
label did more to end inquiry than any amount of evidence ever could have.

None of which means the mundane explanation is usually wrong. It is usually right.
That is the whole point of Occam's razor and I stand by it, with the caveats from
Section~\ref{sec:burdenproof} firmly attached. The overwhelming majority of what gets reported is
aircraft, satellites, balloons, Venus, and honest misperception by decent people who
were not lying about anything. I would put the number well above ninety per cent and
I would not fight you if you said ninety-five.

But ninety-five per cent is not one hundred per cent, and the entire subject lives in
the remainder. The residue. The cases that survive competent investigation by people
who were actively trying to explain them away and could not. That residue is small,
unglamorous, and it is the only thing in Ufology worth a scientist's afternoon.
Doubt driven scepticism cannot see the residue, because it stops working the moment a
plausible-sounding explanation appears, whether or not that explanation was ever
checked.

So doubt everything, including this book. Especially this book. But doubt the way a
mechanic doubts a strange noise --- by opening the hood --- and not the way a man
doubts a knock at the door he would rather not answer.

% ---------------------------------------------------------------
\section{Deconstructing the Peer Review}\label{sec:deconstructingpeerreview}

The last tool in this chapter is an experiment, and it is the only one in the toolkit that
was run on the institutions themselves rather than on a witness. Between 2017 and 2018
Peter Boghossian, James Lindsay and Helen Pluckrose spent ten months writing academic
papers that were false on purpose, engineering the language of each one until it could
clear peer review at respected journals, and then submitting them to find out where the
screening actually failed.\footnote{Boghossian, Lindsay and Pluckrose's own methods write-up
and full paper-by-paper results, together with their long-form account of the project on
\emph{The Joe Rogan Experience}, are in the References under ``New Media.''} Twenty papers
went out. Seven were accepted. That is the whole finding, and everything worth taking from
it is in how those seven were built.

The method was deliberate and it is the part to memorise, because it is the part that
transfers. They did not write a bad paper and hope. They started from a conclusion that was
absurd or indefensible, then went into the existing literature and assembled a chain of
real, correctly formatted citations that made the absurd conclusion read as the next
obvious step. Their word for what they were producing was sophistry --- a forgery of
knowledge, indistinguishable in surface texture from the real thing. The only difference
between their papers and the surrounding literature, as they put it, was that they knew
they had made theirs up.

The accepted papers are worth naming, because the abstract point lands harder when you
hear what actually got through. A paper reporting that dog parks are sites of rape culture
among canines, on the strength of statistics no reviewer could have checked, was published
in a leading journal of feminist geography --- and then singled out by that journal for
special recognition as one of twelve leading pieces in the field for its twenty-fifth
anniversary. A chapter of Hitler's \emph{Mein Kampf}, rewritten in the vocabulary of
solidarity feminism, was accepted by a feminist social-work journal. A paper arguing that
fat bodybuilding should be recognised as a competitive category, on the premise that a fat
body is a legitimately built body, was published in a fat studies journal. Reviewers were
enthusiastic. Four of them liked the work well enough to invite the authors to peer review
other people's papers.

Read that list again with an engineer's eye rather than a partisan's. Nothing in it turns on
whether the reader thinks these fields are worthwhile, and I am not asking anyone to hold an
opinion about them; the fields are the test rig, not the finding. The finding is that the
filter was passing on surface features. Correct vocabulary, correct citation density, correct
posture, conclusions the reviewer already agreed with. Peer review was checking whether the
paper sounded like the field. It was not checking, and structurally could not check, whether
any of it was true. When the authors got better at engineering the surface, their rate of
getting past desk rejection went \emph{up} --- from nothing at the start to almost everything
by the end. The filter did not merely fail to catch the fabrications. It rewarded the ones
built most fluently.

And that is the whole lesson, so let me put it plainly. Peer review is not a truth
detector. It never was. It is a small number of busy people, usually unpaid, usually
anonymous, reading a manuscript for a few hours and deciding whether it looks like
competent work of the sort their field produces. It catches obvious errors and
missing controls. It does not replicate your experiment, audit your data, or visit
your lab. It is a plausibility filter, and plausibility is exactly the thing a
motivated author can manufacture.

Which cuts two ways for us, and I want to be very careful here, because this section
is the easiest one in the book to misuse.

The first way is a warning against a bad habit in my own community. ``Peer review is
broken'' has become the universal solvent of the fringe --- a reason to dismiss any
inconvenient finding while accepting anything published on a blog. That is doubt
driven scepticism from Section~\ref{sec:doubtdrivenskepticism} with a lab coat on. Boghossian, Lindsay and
Pluckrose did not show that peer review is worthless. They showed it has a specific,
diagnosable failure mode: it fails where a field's conclusions are settled in advance
and the review process is checking for conformity rather than for error. A journal of
solid-state physics is not vulnerable in the same way, because the referee can check
your numbers against a lattice constant that does not care about anyone's politics.

The second way is the one that concerns Ufology directly, and it is the harder
lesson. If a field can be captured by the assumptions of its practitioners, then this
field is captured to the bone. Consider what our review process actually looks like.
A case is investigated by an organisation whose members joined because they find the
subject compelling. It is written up for a publication read almost exclusively by
people who already accept the premise. It is refereed, when it is refereed at all, by
colleagues who share the founding assumption. Then it is cited by the next
investigator as established. That is not peer review. That is a field checking
whether new work sounds like the work it already produced --- which is precisely the
failure the grievance studies papers walked through.

And if you want the cheap version of the fix, it is the casting sheet from the opening pages of
this chapter. Put a competent hostile reader in the room and pay him to stay. \emph{UFO~Hunters}
had no funding, no laboratory and no institutional standing, and it still produced better
evidential hygiene in season one than most of the published literature, purely because the sceptic
was contractually present and could not be edited out of the argument. When his chair emptied, the
standard fell. That is the whole of peer review in one sentence, and it is available to anyone
willing to be embarrassed in front of a witness.

We do it in the other direction too. The debunking literature has its own
orthodoxies, its own approved vocabulary, its own conclusions known before the file is
opened. A paper concluding ``inconclusive, requires further instrumentation'' is
unpublishable in either camp. Both audiences want a verdict.

So what do we do about it? Not much, honestly, until the field grows up --- but here
is what I try to hold myself to, and what I would ask of you.

\begin{itemize}[leftmargin=1.4em,itemsep=3pt]
  \item \textbf{Prefer review from outside the tent.} An anomalous metal sample should
  go to a metallurgist who has never heard of Roswell and has no stake in the answer.
  If the finding is real it will survive a hostile lab. If it needs a friendly one, it
  is not a finding.
  \item \textbf{Publish the failures.} Section~\ref{sec:claimsmadeslipperyslope} faulted \emph{Ancient Aliens} for
  editing out the tests that went badly. We do the same thing when we quietly drop the
  cases that resolved. The resolved cases are how you calibrate. They are the
  denominator, and a field with no denominator can prove anything.
  \item \textbf{Cite first appearances, not repetitions.} Section~\ref{sec:confirmationbias} covered this.
  Peer review does not catch source laundering; only you can.
  \item \textbf{State your priors in writing.} If a reader knows what you believed
  before you began, they can weight your conclusion appropriately. This is not a
  confession. It is data about the instrument.
\end{itemize}

I am aware of the irony of this section. Nothing in this book has been peer reviewed.
I am a self-appointed professor of a field that does not grant degrees, publishing my
lecture notes and asking you to take them seriously. So do not take them seriously ---
take them provisionally. Check the sources at the back. Trace the claims. Where I have
reasoned past the edge of the evidence, say so out loud, because a visible hole in an argument is
worth more than a seamless one you cannot inspect.

That is the end of the toolkit. Language distorts the question before you ask it. The
burden of proof sits with whoever makes the claim, including me. Evidence is not proof,
and proof may never come. One good idea does not license a hundred bad ones.
Correlation is not causation. Your brain will curate the evidence on your behalf
without asking permission. Doubt is a method and not a destination. And the
institutions we trust to check all of this are checking something rather less than we
assume.

Keep all seven of those within reach. Starting in the next chapter, we are going to
need every one of them.

And keep the three chairs in mind as well, because the seven tools are what the chairs are for.
Birnes will tell you what the evidence could mean, and you need him, because a field with no
believer never goes and looks. Acworth will tell you what the evidence actually measures, and you
need him more, because a field with no sceptic never finds out. Uskert sits between them and
actually moves when the evidence moves, which is the only honest place for the rest of us to stand.
\hi{You will not always have a television budget and a lab, but you can always cast all three parts
in your own head --- and the only unforgivable move is to let one of them do all the talking.}

\begin{srcnotes}
\item \emph{UFO~Hunters} (History Channel, 2008--2009), three seasons, produced by Motion Picture
Productions for A\&E Television Networks: the investigative format built around one believer
(William~J.\ Birnes), one sceptic with laboratory training (Ted Acworth, seasons one and two), and
one undecided investigator (Kevin Cook, then Pat Uskert), used as the frame for this chapter.
\item Birnes, W.\ J.\ \emph{UFO Magazine} (publisher and editor), and Birnes, W.\ J.\ \& Corso,
P.\ J.\ \emph{The Day After Roswell} (Pocket Books, 1997), for the evidential standards of the
believer seat discussed in Section~\ref{sec:confirmationbias}.
\item Acworth, T., biographical and professional record: MIT and Stanford mechanical engineering,
founder of Artaic; his on-camera instrumentation and measurement work in
\emph{UFO~Hunters} seasons one and two is the basis of the falsifiability argument in
Section~\ref{sec:burdenproof}.
\item Uskert, P., \emph{UFO Road Trip} and his contributions across all three seasons of
\emph{UFO~Hunters}, for the undecided position described in
Section~\ref{sec:doubtdrivenskepticism}.
\item Wason, P.\ C.\ (1960). On the failure to eliminate hypotheses in a conceptual task.
\emph{Quarterly Journal of Experimental Psychology}, 12(3), 129--140 --- the 2--4--6 task quoted
in Section~\ref{sec:confirmationbias}.
\item Popper, K.\ R.\ \emph{The Logic of Scientific Discovery} (Hutchinson, 1959) and
\emph{Conjectures and Refutations} (Routledge, 1963), for falsifiability as used throughout
Section~\ref{sec:burdenproof}.
\item Sagan, C.\ \emph{The Demon-Haunted World: Science as a Candle in the Dark} (Random House,
1995), for the baloney-detection framing behind this chapter's toolkit.
\item Boghossian, P., Lindsay, J.\ \& Pluckrose, H.\ (2017--2018). The grievance studies affair:
twenty deliberately fabricated papers submitted to peer reviewed journals, seven accepted, as
discussed in Section~\ref{sec:deconstructingpeerreview}. The named cases in that section are the
dog park paper in \emph{Gender, Place and Culture} (later selected for the journal's twenty-fifth
anniversary recognition), the \emph{Mein Kampf} rewrite accepted by \emph{Affilia}, and the fat
bodybuilding paper in \emph{Fat Studies}. See also the earlier retracted \emph{Cogent Social
Sciences} paper (Boghossian \& Lindsay, 2017). Authors' own write-up and the long-form interview
are listed in the References under ``New Media.''
\end{srcnotes}
